\chapter{Configuración del entorno y datos}

\section{Fuentes de datos}

El sistema trabaja con varias fuentes porque el mercado de CS2 está fragmentado. Ninguna fuente aislada contiene toda la información necesaria para decidir. Steam aporta precio, histórico y señales del mercado principal; BUFF163 permite comparar oportunidades y buy orders; SteamDT ayuda a descubrir candidatos; el OCR permite capturar datos cuando no existe una API cómoda; y los CSV o importaciones manuales sirven para integrar muestras controladas o material histórico.

\subsection{Mercado de Steam}

Steam actúa como una de las fuentes principales de precio e histórico. En el proyecto se capturan variantes reales de artículos, precios actuales, moneda, buy orders y puntos históricos cuando están disponibles. La información se normaliza para que pueda combinarse con otras plataformas y para evitar que las decisiones dependan de nombres o formatos propios de una fuente.

\subsection{Mercados externos}

BUFF163 se utiliza como mercado externo de comparación. Sus precios se reciben principalmente en CNY, por lo que el sistema conserva tanto el valor original como su conversión a EUR. Esta doble representación permite auditar el dato original y, al mismo tiempo, entrenar modelos sobre una moneda canónica. SteamDT se utiliza como apoyo para descubrir candidatos y estrategias de búsqueda, no como fuente única de decisión.

\section{Adquisición y persistencia}

\subsection{Pipeline de adquisición}

La adquisición se organiza como una capa de agentes extractores. Estos procesos pueden ejecutarse como comandos CLI, jobs o servicios programados. Su función es capturar datos, normalizarlos, validar errores y guardarlos con trazabilidad. El flujo incluye scraping, refresco de históricos, OCR, importación manual y un bucle operativo local que puede ejecutar scraping y actualización periódica.

La adquisición periódica contempla cookies o sesiones, delays configurables, límites de concurrencia, reintentos, deduplicación y registro de anomalías. Esta parte es importante porque las fuentes no siempre son estables: pueden cambiar el formato, limitar el ritmo de acceso o devolver información incompleta.

\subsection{Modelo de persistencia}

La persistencia se apoya en Supabase/Postgres. La decisión de diseño es tratar la base de datos como fuente de verdad y mantener el histórico de forma append-only siempre que sea posible. Las tablas operativas principales son \texttt{market\_items} y \texttt{market\_history\_points}. La primera mantiene una fila por variante real de artículo y estado actual por plataforma; la segunda guarda series temporales en formato largo por artículo, plataforma, fecha y métrica.

También se define un esquema canónico más amplio con entidades como \texttt{assets}, \texttt{platforms}, \texttt{market\_observations}, \texttt{predictions}, \texttt{investment\_decisions} y \texttt{simulated\_positions}. Esta separación permite evolucionar desde el esquema operativo actual hacia una representación más trazable y auditable.

\section{Preparación de características}

\subsection{Variables de mercado}

Las características de mercado incluyen precios Steam y BUFF normalizados, precios originales, spread bruto, spread neto, buy orders, volumen, liquidez, número de listings, antigüedad de variante y señales temporales. En la fase supervisada se exploraron variables de momentum y reversión como retornos a 1, 3 y 7 días, distancia a medias móviles, RSI, volatilidad y desviaciones de ventas.

La exploración inicial detectó que algunas señales temporales aportan información moderada, mientras que ciertas variables categóricas de alta cardinalidad, como identificadores concretos de skins, deben tratarse con cuidado para evitar sobreajuste. Por ello, el entrenamiento separa variables de evaluación y variables de entrenamiento, y permite realizar ablations sin reconstruir todo el dataset.

\subsection{Variables de cartera}

Las variables de cartera proceden del simulador. Incluyen saldo disponible, valor de posiciones abiertas, capital bloqueado por \textit{trade hold}, fecha de desbloqueo, exposición por activo, exposición por plataforma, beneficio realizado y drawdown. Estas variables son necesarias para que el agente Portfolio no decida únicamente por oportunidad local, sino por estado global de la cartera.

\section{Simulador de cartera}

El simulador reproduce las reglas económicas que antes estaban dispersas en el Excel operativo. No intenta copiar la hoja de cálculo, sino migrar sus reglas útiles a código comprobable: conversión EUR/CNY, comisiones, beneficio neto, rentabilidad porcentual, fecha de desbloqueo y capital bloqueado.

\subsection{Operaciones permitidas}

Las operaciones básicas son comprar, vender y mantener. Una compra abre una posición con precio de entrada, cantidad, plataforma y fecha de bloqueo. Una venta cierra una posición si está disponible y calcula el beneficio realizado. Mantener no modifica la posición, pero puede afectar al rendimiento si el mercado se mueve o si el capital sigue inmovilizado.

\subsection{Comisiones y valoración}

El Excel operativo utiliza reglas de referencia que se han convertido en hipótesis de simulación. BUFF se modela con un factor neto de venta de 0,975. Steam usa una comisión de mercado del 0,87 y, en escenarios conservadores de salida a banco, puede aplicarse una pérdida adicional del 20\%. La conversión inicial simplificada usa 1 EUR = 8 CNY, aunque queda pendiente versionar tipos de cambio por fecha si se incorporan fuentes externas.

El \textit{trade hold} se modela de forma conservadora como ocho días desde la compra, equivalente a siete días más una ventana de desbloqueo. Esta simplificación evita declarar vendible una posición demasiado pronto.

\section{Interfaz de consulta}

El proyecto incluye una interfaz web como apoyo operativo. Su objetivo no es vender el sistema como producto final, sino hacer visible el estado de adquisición, modelos, simulador, restricciones y recomendaciones. La interfaz debe permitir responder rápidamente qué datos existen, qué modelo está activo, qué señales se están generando y por qué una oportunidad se recomienda, se observa o se descarta.
